\documentclass[11pt,openany]{book}
\usepackage[margin=1in]{geometry}
\usepackage{amsmath,amssymb,amsthm}
\usepackage{graphicx}
\usepackage{booktabs}
\usepackage{longtable}
\usepackage{listings}
\usepackage{hyperref}
\usepackage{xcolor}
\usepackage{enumitem}
\usepackage{tikz}
\usetikzlibrary{arrows.meta,positioning,fit,calc}

\definecolor{cb-bg}{RGB}{245,245,245}
\definecolor{cb-kw}{RGB}{130,30,130}
\definecolor{cb-str}{RGB}{30,120,30}
\definecolor{cb-com}{RGB}{120,120,120}
\lstset{
  basicstyle=\ttfamily\footnotesize,
  backgroundcolor=\color{cb-bg},
  keywordstyle=\color{cb-kw}\bfseries,
  stringstyle=\color{cb-str},
  commentstyle=\color{cb-com}\itshape,
  breaklines=true,
  frame=single,
  framesep=4pt,
  showstringspaces=false,
  captionpos=b,
}
\lstdefinelanguage{pyshell}{
  keywords={def,class,return,if,else,elif,for,while,with,import,from,as,yield,raise,try,except,pass,None,True,False,lambda,self,async,await},
  sensitive=true,
  comment=[l]{\#},
  morestring=[b]",
  morestring=[b]',
}

\theoremstyle{definition}
\newtheorem{definition}{Definition}[chapter]
\newtheorem{observation}{Observation}[chapter]
\newtheorem{conjecture}{Conjecture}[chapter]

\theoremstyle{plain}
\newtheorem{result}{Result}[chapter]

\newcommand{\bpb}{\text{bpb}}
\newcommand{\tokpb}{t_{\text{per byte}}}
\newcommand{\toks}{\text{tok}/\text{s}}
\newcommand{\effdim}{\mathrm{EffDim}}
\newcommand{\conr}{\mathrm{ConR^2}}
\newcommand{\posr}{\mathrm{PosR^2}}
\newcommand{\cka}{\mathrm{CKA}}
\newcommand{\omegaproj}{W_\omega}

\title{%
  \textbf{Chronohorn} \\
  \large The Carving Machine: \\
  Byte-Level Language Modeling, \\
  Substrate Dynamics, and the Frontier \\[1em]
  \normalsize Sessions 6--11, Architecture, Forensics, and Infrastructure}
\author{asuramaya \\ with Claude Opus 4.7}
\date{April 18, 2026}

\begin{document}
\frontmatter
\maketitle

\tableofcontents

\chapter*{Foreword}
\addcontentsline{toc}{chapter}{Foreword}

This book is the accumulated state of a research program called \emph{Chronohorn}:
an attempt to find a language-model family that achieves one bit per byte at one
million tokens per second on a single consumer GPU --- specifically, the
NVIDIA RTX A4000, a 16\,GB card whose memory bandwidth and compute envelope
are representative of what an individual researcher, rather than a corporate
cluster, can keep warm. As of this writing we have 1.785 bpb on one
architecture and 324{,}000 tok/s on another --- an $8\times$ speed gap and a
$0.8$ bpb quality gap from the goal. The gap has been halved every session
since Session 6. The path forward is visible but non-obvious, and most of
this book is a record of the non-obvious choices that narrowed it.

The book has five parts. Part I establishes the problem: what one bit per byte
means, how we measure it without cheating, and why the naive choices
($2000$-batch evaluation, 200-parameter probes, sentencepiece tokenizers)
get in the way of the answer. Part II is a catalog of architectural primitives
--- substrates, rotations, readouts --- with enough mathematical grounding
that anyone who finds one of them confusing can locate the confusion in a
specific equation. Part III is a chronicle: Session 6 through Session 11, in
the order we actually learned the things. Part IV is about forensics, the
tools we use to open a model that has reached a plateau and understand why it
plateaued. Part V is about the infrastructure --- the fleet, the runtime
daemon, the database, the MCP surface, the dashboard --- that allows a single
researcher to keep dozens of experiments in flight and never lose a result
again. Appendices cover command-line reference, database schema, and file
formats.

A note on the voice. This is a technical book but not a paper. It is written
the way the project is conducted: opinionated, specific, with named
mechanisms that failed and named mechanisms that worked. When we call a
primitive ``dead'' we mean it was ruled out with evidence, and Chapter
\ref{ch:forensics} will explain exactly what the evidence was. When we call
something ``the frontier'' we mean we have a number and a date. The goal is
that a future researcher --- human or model --- can walk into this project
cold, read Part III, and know not just what the best result is but which
dead ends have already been walked.

\mainmatter

% ============================================================
\part{Foundation}
% ============================================================

\chapter{The Problem}
\label{ch:problem}

\section{One Bit per Byte, One Million Tokens per Second}

A language model, at the layer of measurement we care about, is a
conditional probability distribution over the next byte of text given some
prefix. Call the model $p_\theta(x_t \mid x_{<t})$. Given a test corpus of
length $N$ bytes, the \emph{bits per byte} is
\begin{equation}
  \bpb = -\frac{1}{N}\sum_{t=1}^N \log_2 p_\theta(x_t \mid x_{<t}).
  \label{eq:bpb}
\end{equation}

This is not a proxy. Equation \ref{eq:bpb} is, up to the choice of $p_\theta$,
the quantity that determines how much a well-designed compression codec can
squeeze the text down to. On English, the Shannon-attributed lower bound is
around $0.6{-}1.3$ bits per character depending on the corpus; on modern
wide-domain datasets like FineWeb it sits around $0.7{-}0.9$ for state of
the art and $1.0$ is a widely recognized round-number threshold for
``good''.

The speed target is $10^6$ tokens per second, measured on a single A4000. For
byte-level models, a token \emph{is} a byte; for subword models (sp1024 in
our ecosystem), a token is $\sim 2.4$ bytes. Speed will be discussed in
Chapter \ref{ch:speed} once the architecture has been named; for now it
suffices to note that the target is unforgiving: the A4000's HBM bandwidth
tops out around 448\,GB/s, and its tensor core throughput is about 19.2
TFLOPs at FP16. At $10^6$ tok/s, you have $19.2$ TFLOPs per million bytes,
or $\sim 19$ MFLOPs per byte --- a number that rules out most architectures
whose forward pass is dominated by dense $d \times d$ matmuls with $d \geq
512$.

\section{Bits per Byte versus Bits per Token}

A subtle and repeatedly rediscovered mistake is to report the cross-entropy
in the token stream the model was trained on. If the model uses a
sentencepiece tokenizer that, on the test corpus, averages 2.436 bytes per
token, then
\begin{equation}
  \bpb = \bpb_{\text{token}} \cdot \tokpb = \frac{\bpb_{\text{token}}}{\text{bytes per token}}
  \label{eq:bpb-conversion}
\end{equation}
where $\tokpb = 1 / (\text{bytes per token}) \approx 0.411$ for sp1024 on
FineWeb 10B. A model reporting 1.80 ``bpb'' on its own tokenizer is really
reporting $1.80 \times 2.436 = 4.38$ bpb in bytes, which is much worse than
a decent byte-level model.

Chronohorn enforces Equation \ref{eq:bpb-conversion} at ingestion: the
\texttt{record\_result} function computes $\bpb$ from
\texttt{model.test\_bits\_per\_token} using $\tokpb$ from the dataset
section of the result JSON, and rejects results without the conversion
factor. Every number reported after Chapter \ref{ch:problem} of this book,
unless otherwise annotated, is byte-normalized.

\section{Probes versus Finals}

A second mistake we repeat is to compare probes (short evaluations used to
monitor training) across experiments. Probes at Chronohorn are 2--8
evaluation batches; a final is 32--200 batches depending on configuration.
At 2 batches, the standard deviation of the bpb estimate is large enough
(observed $\sim 0.03$) that two probes can differ by 0.05 bpb on identical
models. At 32 batches the standard deviation is small enough ($\sim 0.005$)
that a 0.02 bpb difference is meaningful.

Chronohorn's probe rows carry an \texttt{eval\_batches} column in the database.
The dashboard displays probe bpb with a tilde (``$\sim$'') suffix and only
compares finals across experiments. Comparing a probe to a final is
treated the same way as comparing a token-level bpb to a byte-level bpb: a
type error.

\section{Causality}

A byte-level language model must be strictly causal: the output at position
$t$ must depend only on inputs at positions $\leq t$. A causality violation
at position $t$ --- a leak from $t+1$ into $t$ --- can silently improve bpb
by a small and seductive amount, making the model appear better than it
is and producing generation that is incoherent because the prefix contains
information the generator doesn't yet have.

We enforce causality with \texttt{decepticons/tests/test\_causality.py}: the
test feeds two sequences identical up to position $t$ and different after
$t$, then asserts the logits at $t$ are bit-identical. Every architectural
primitive must pass this test before it can be added to the substrate.

\section{Why Bytes, Not Tokens}

Three reasons:
\begin{enumerate}[label=\roman*.]
  \item The tokenizer is a frozen, black-box preprocessor. If the
    architecture learns representations that are interesting only in
    token-space, they don't transfer to byte streams, which is what actual
    text, code, JSON, and network protocols are.
  \item A 1024-token vocabulary hides the difficulty of the local problem.
    The model rarely has to predict ``\verb|o|'' from ``\verb|hell|'', because
    ``\verb|hello|'' is often one token. A byte-level model sees each
    completion and has nowhere to hide.
  \item Tokenizer boundaries leak. Sentencepiece splits on whitespace in a
    way that correlates with content but is not determined by the model; this
    creates spurious structure that can absorb useful gradient.
\end{enumerate}

The compounding cost is that a byte-level model must do roughly $2.4\times$
more token-steps per byte of text. This is what makes the $10^6$ tok/s
target hard.

% ------------------------------------------------------------
\chapter{The Ecosystem}
\label{ch:ecosystem}

\section{Three Codebases}

Chronohorn is one of three codebases the project spans. The split was
introduced during Session 7 to keep the \emph{what} of the model separate
from the \emph{how} of running experiments.

\begin{center}
\begin{tikzpicture}[node distance=1.2cm and 2.8cm,
  box/.style={rectangle,draw,rounded corners=3pt,minimum width=3.2cm,minimum height=1cm,align=center,font=\small},
  arrow/.style={->,>=Stealth,thick}]
\node[box] (dec) {\textbf{decepticons}\\ \footnotesize kernels, substrates,\\ \footnotesize readouts};
\node[box, right=of dec] (ch)  {\textbf{chronohorn}\\ \footnotesize runtime, tracker,\\ \footnotesize fleet, MCP};
\node[box, right=of ch]  (he)  {\textbf{heinrich}\\ \footnotesize forensics, MRI,\\ \footnotesize weight analysis};
\draw[arrow] (ch) -- (dec) node[midway, above] {\footnotesize imports};
\draw[arrow] (he) -- (ch)  node[midway, above] {\footnotesize reads DB};
\draw[arrow] (he) to [bend right=30] node[below] {\footnotesize reads checkpoints} (dec);
\end{tikzpicture}
\end{center}

\begin{itemize}
  \item \textbf{decepticons} is the model library. It defines the
    causal-bank substrate, rotation primitives, readout kinds, adaptive
    substrate transforms, and the Triton scan kernel. It has no knowledge of
    experiments, databases, or the fleet. It can be used standalone by a
    researcher writing a training script by hand.
  \item \textbf{chronohorn} is the runtime. It owns the SQLite database
    of experiments, the fleet-dispatch logic, the HTTP dashboard, the MCP
    tool surface (55 tools), and the drain loop that pulls results back
    from k8s pods. It depends on decepticons but decepticons does not
    depend on it.
  \item \textbf{heinrich} is a forensic suite. Given a checkpoint, it
    computes effective dimensionality, content and position regression
    $R^2$, CKA between layers, routing-entropy histories, and the singular
    value spectrum of $\omegaproj$. It reads checkpoints from the shared
    Sharts volume and writes MRI artifacts back.
\end{itemize}

The dependency is unidirectional: \texttt{chronohorn $\to$ decepticons},
never the reverse. We enforce this with a ``firewall test'' in the
decepticons CI that grep-rejects any \texttt{import chronohorn} line.
Session 8 restored this invariant after a slow accumulation of cross-imports
had made the two codebases structurally coupled; the split now stands.

\section{The Fleet}

Chronohorn dispatches experiments to a small k3s cluster of three hosts:
\texttt{slop-01}, \texttt{slop-02}, and \texttt{slop-home}. slop-01 is the
k3s server and LAN gateway (192.168.4.129); slop-02 and slop-home are
agents. Each host has one GPU:
\begin{itemize}
  \item \texttt{slop-01} and \texttt{slop-02}: RTX A4000 (Ampere, sm\_86),
    16\,GB, $\sim$19 TFLOPs FP16.
  \item \texttt{slop-home}: Quadro RTX 4000 (Turing, sm\_75), 8\,GB,
    $\sim$7 TFLOPs FP16. Known limitation: \texttt{pytorch.compile} with
    reduce-overhead kernels triggers a device-side assert via Inductor's
    Triton autotune on Turing. We skip slop-home for byte-level runs.
\end{itemize}

Experiments are shipped as k8s Jobs with a managed-command spec that:
(i) installs system deps (gcc, sentencepiece), (ii) snapshots chronohorn's
\texttt{python/} and decepticons's \texttt{src/} into \texttt{/run/source},
(iii) runs the trainer under \texttt{PYTHONDONTWRITEBYTECODE} with results
written to \texttt{/run/results/}, which is a hostPath mount to
\texttt{/tmp/chronohorn-runs/<job>/results/} on the node.

% ------------------------------------------------------------
\chapter{Measurement Methodology}
\label{ch:methodology}

\section{The Eval Stream}

FineWeb 10B in byte form occupies $\sim$10\,GB of raw text; our test split
is a fixed $\sim$100\,MB slice. During training, we evaluate the model on
the test split at probe checkpoints. The stream is reset to byte 0 before
every probe, guaranteeing that every probe and final measures the same
data. Without this reset, a probe at step 1000 sees different bytes than a
probe at step 2000, and neither is the same as a probe at step 20000 ---
making probe-to-probe trajectories uninterpretable.

\begin{lstlisting}[language=pyshell,caption={Eval stream reset (simplified)}]
def evaluate(model, test_loader, max_batches):
    test_loader.reset_to_position(0)  # deterministic
    total_loss = 0.0
    total_tokens = 0
    for i, (x, y) in enumerate(test_loader):
        if i >= max_batches:
            break
        with torch.no_grad():
            logits = model(x)
            loss = F.cross_entropy(logits.view(-1, V), y.view(-1))
        total_loss += loss.item() * x.numel()
        total_tokens += x.numel()
    return total_loss / total_tokens  # nats; convert later
\end{lstlisting}

\section{Probe Policy}

The probe policy balances cost against signal. At step 0, training has zero
information; the probe loss is $\approx \log 256$ nats / byte, which is 8.0
bpb. Probe cost is dominated by the forward pass over \texttt{max\_batches}
batches of sequence length $L$; at $L=1024$, $B=32$, $M=16$, a probe is
$524$k tokens.

The \emph{adaptive} probe policy (default since Session 8):
\begin{itemize}
  \item Probe at a geometric sequence: step 200, 400, 800, 1600, 3200, 6400.
  \item \emph{Micro} probes (4 batches, 128k tokens) before step 800.
  \item \emph{Standard} probes (8 batches, 256k tokens) after step 800.
  \item \emph{Promotion} probes (16 batches, 524k tokens) at every
    checkpoint-save step, every $5000$ steps.
\end{itemize}

The micro-to-standard promotion is controlled by \texttt{--probe-promotion-count}:
after $N$ standard probes in a row, a promotion probe fires. Promotion probes
are what land in the Chronohorn frontier table as the ``final'' when the
training script exits at 20k steps.

A 32-batch final is $1.05$M tokens and takes $\sim$3 seconds on an A4000;
a 200-batch final is $\sim$20 seconds. We budget both.

\section{Comparing Across Runs}

Three rules, enforced:
\begin{enumerate}
  \item Same eval batch count (usually 32) for any bpb comparison.
  \item Same tokenizer for any bpb comparison (or use
    Equation \ref{eq:bpb-conversion}).
  \item Same test stream position reset policy.
\end{enumerate}

The dashboard's frontier table and the MCP \texttt{chronohorn\_compare} tool
both enforce rule 1 by filtering on \texttt{eval\_batches}. Rule 2 is
enforced by the family adapter's \texttt{infer\_from\_config} path, which
rejects comparisons across families without an explicit override.

% ============================================================
\part{The Architecture Catalog}
% ============================================================

\chapter{Causal-Bank: The Frozen Substrate}
\label{ch:causal-bank}

The causal-bank family is the mainline architecture of this book. It was
introduced in Session 6 (documented in \texttt{paper/session6\_chronohorn.tex})
and has been mutated, probed, and scaled continuously since then. This
chapter establishes the base model before introducing the mechanisms that
have been layered on top of it.

\section{Anatomy of a Causal Bank}

A causal-bank model in forward order:
\begin{enumerate}
  \item \textbf{Embedding.} Input bytes $x_t \in \{0,\ldots,255\}$ are
    embedded into $\mathbb{R}^d$ with a learned $256 \times d$ matrix. In
    our frontier configurations $d = 320$; at scale 12, $d = 480$.
  \item \textbf{Local path.} A small causal convolution with kernel
    $K=8$ applies a learned linear transform over the last 8 bytes of
    context. Output $\in \mathbb{R}^d$. Can be disabled with
    \texttt{--local-scale-override 0.0}.
  \item \textbf{Substrate.} A linear dynamical system with $M$ modes. At
    time $t$, the substrate state $h_t \in \mathbb{R}^M$ is
    \begin{equation}
      h_t = \alpha \odot h_{t-1} + (1 - \alpha) \odot W_{\text{in}} x_t,
      \label{eq:substrate}
    \end{equation}
    where $\alpha \in [0,1]^M$ is a per-mode decay (EMA factor), $\odot$
    is elementwise multiplication, and $W_{\text{in}} \in
    \mathbb{R}^{M \times d}$ projects the embedded input into mode space.
    In the \emph{frozen} substrate, $W_{\text{in}}$ is fixed at
    initialization; $\alpha$ is parameterized by a half-life
    distribution over $[\tau_{\min}, \tau_{\max}]$.
  \item \textbf{Readout.} A function $f_\theta: \mathbb{R}^M \to
    \mathbb{R}^V$ (where $V = 256$ for byte vocab) that maps substrate
    state to logits. This is where most of the free parameters live;
    Chapter \ref{ch:readouts} covers the variants.
\end{enumerate}

The loss is cross-entropy on the next byte:
\begin{equation}
  \mathcal{L} = -\mathbb{E}_t\left[\log p_\theta(x_{t+1} \mid x_{\leq t})\right],
  \quad p_\theta = \mathrm{softmax}(f_\theta(h_t)).
\end{equation}

\section{Half-Life Distribution}

The per-mode decay $\alpha_i$ is parameterized as
\begin{equation}
  \alpha_i = \exp\left(-\frac{\ln 2}{\tau_i}\right),
  \quad \tau_i \in [\tau_{\min}, \tau_{\max}]
\end{equation}
where $\tau_i$ is the \emph{half-life} of mode $i$ in tokens. A mode with
$\tau = 8$ decays to half its value in 8 tokens; a mode with $\tau = 512$
takes 512. In a byte-level model at $L = 1024$, a mode with $\tau = 2048$
has barely decayed at all by the end of the sequence --- that's a
long-context mode.

The \texttt{--linear-half-life-max} flag controls $\tau_{\max}$. Default
is 512. At scale 8 we typically use $\tau_{\min} = 2$, $\tau_{\max} =
512$, distributed log-uniformly across modes.

\section{Frozen versus Adaptive}

\begin{definition}[Frozen substrate]
The substrate is \emph{frozen} if $W_{\text{in}}$ in Equation
\ref{eq:substrate} is fixed at initialization and not trained.
\end{definition}

\begin{definition}[Adaptive substrate]
The substrate is \emph{adaptive} if there exists a learned projection
$\omegaproj: \mathbb{R}^d \to \mathbb{R}^M$ that produces an input-dependent
modulation $\omega_t = \omegaproj x_t$, and the substrate update is
\begin{equation}
  h_t = \alpha \odot h_{t-1} + (1 - \alpha) \odot f_{\text{HRR}}(W_{\text{in}} x_t, \omega_t)
\end{equation}
where $f_{\text{HRR}}$ is a Holographic Reduced Representation mixing
operation (typically elementwise multiplication in a Fourier basis).
\end{definition}

The adaptive substrate is controlled by \texttt{--adaptive-substrate}.
The flag \texttt{--hrr-omega-init} initializes $\omegaproj$'s bias to a
Fourier basis (one frequency per mode) and the weight to \emph{zero}.
As we will see in Chapter \ref{ch:forensics-omega}, this zero-initialization
is load-bearing and creates a structural gradient block in some readout
configurations.

\section{Params at Scale}

The parameter count at scale $s$ is dominated by the readout. At $s=8$,
$d=320$, $M=2048$ (the defaults), a model with routed-experts readout
has $\sim 18$\,M parameters; an MLP readout has $\sim 17.8$\,M; a
tied-embedding readout has $\sim 5$\,M. These are the scales at which
Session 9--11 experiments were conducted.

% ------------------------------------------------------------
\chapter{Rotation Primitives}
\label{ch:rotations}

Rotations are the Session 9 innovation: input-dependent, noncommutative
transformations applied during the substrate scan. The goal of a rotation
is to carry \emph{order-sensitive} information through the EMA without
overwriting it.

\section{Why Rotations at All}

The bare EMA substrate (Equation \ref{eq:substrate}) is a scalar-per-mode
recurrence: mode $i$ at time $t$ depends on mode $i$ at time $t-1$ and the
current input. It cannot express a function of the form ``A happened
before B''; the EMA mixes contributions order-invariantly within each
mode.

A rotation adds a matrix-per-mode or matrix-per-pair recurrence: the state
is rotated by an input-dependent operator before the EMA mixing. If the
operator group is non-abelian, then $A \circ B \neq B \circ A$ and the
network can express order-sensitivity.

Session 9 showed (Chapter \ref{ch:session9}) that order is not the
bottleneck at the bpb range we're in --- $\posr \approx 0.001$ across all
configurations --- but that the rotation machinery improves
\emph{content} coupling. The lasso rotation doubled $\conr$ from 0.101 to
0.209, and this delivered $-0.109$ bpb.

\section{The Lasso}
\label{sec:lasso}

The lasso rotation is a $2 \times 2$ matrix transition per mode pair. For
modes paired $(2i, 2i+1)$, the state vector $(h_t^{2i}, h_t^{2i+1})^\top$
is rotated by an input-dependent matrix:
\begin{equation}
  \begin{pmatrix} h_t^{2i} \\ h_t^{2i+1} \end{pmatrix}
  =
  R_t^i \begin{pmatrix} h_{t-1}^{2i} \\ h_{t-1}^{2i+1} \end{pmatrix} + \text{(input)}
\end{equation}
where $R_t^i$ is constructed as a general $2 \times 2$ matrix
parameterized by input-dependent entries --- specifically, a learned
projection of the embedding to four real numbers per mode pair.

This is non-abelian: for two input sequences $A$ followed by $B$ vs $B$
followed by $A$, the resulting state differs because $R^B R^A \neq R^A
R^B$ in general.

\begin{result}[Lasso, 2026-04-13]
At scale 8, 20k steps, byte level: $-0.109$\,bpb relative to the
no-rotation baseline at identical capacity. The best frontier model
(legal, 1.842 bpb) uses the lasso.
\end{result}

The scan is implemented via parallel Hillis--Steele prefix scan using
\texttt{torch.roll} and \texttt{torch.where} for dynamo compatibility.
The Triton fused scan (Chapter \ref{ch:speed}) brings throughput to
$\sim$316k tok/s on A4000.

\section{Complex Rotation}

A simpler predecessor to the lasso: an $SO(2)$ rotation with input-dependent
angle.
\begin{equation}
  R_t^i = \begin{pmatrix}
    \cos\theta_t^i & -\sin\theta_t^i \\
    \sin\theta_t^i & \cos\theta_t^i
  \end{pmatrix},
  \quad \theta_t^i = \phi(x_t)_i.
\end{equation}

$SO(2)$ is abelian: $R^A R^B = R^B R^A$. This means complex rotation
\emph{cannot} encode order. Session 9 showed empirically that complex
rotation delivers $-0.008$\,bpb --- basically nothing, consistent with
the prediction that order isn't the bottleneck.

Complex rotation is kept in the codebase as a baseline for rotation
mechanisms; do not use it as a mainline.

\section{Quaternion Rotation}

$SO(3)$ rotation using Hamilton product. Per mode triple $(3i, 3i+1, 3i+2)$,
apply a unit-quaternion rotation to the 3-d state vector. Non-abelian.

\begin{result}[Quaternion, 2026-04-14]
At scale 8, 20k steps: matches lasso in bpb ($-0.107$) but is $7\%$ slower
due to the Hamilton product being more expensive than the $2\times 2$
matrix multiply.
\end{result}

An added unit-norm constraint on the quaternion did not help; the
unconstrained quaternion already converges to near-unit magnitudes during
training.

\section{$SO(5)$ Rotation}

A higher-dimensional rotation using \texttt{matrix\_exp} over a
skew-symmetric Lie-algebra element. Each mode-group of 5 is rotated by
an $SO(5)$ matrix. Requires \texttt{--state-dim 40 --num-heads 8}
(head\_dim = 5).

\begin{result}[SO(5), 2026-04-14]
Stable across all training lengths (the matrix exponential cannot
diverge), but $\sim 2.7\times$ slower than lasso. bpb parity or slightly
worse. Reason: each mode participates in a single 5-d rotation group, so
the total expressiveness per parameter is lower than the lasso's per-pair
expressiveness.
\end{result}

$SO(5)$ is non-solvable and compact, which was the theoretical motivation:
it contains enough structure to express genuinely complex sequential
transformations. Empirically the cost exceeded the benefit.

\section{Quintic Rotation}

A GL(5) unconstrained $5 \times 5$ matrix per mode-group. No normalization.

\begin{result}[Quintic, 2026-04-14]
\textbf{Unstable.} Training diverges to NaN within the first 1000 steps
at most seeds. The unconstrained matrix allows one singular value to
blow up, after which the forward pass saturates. \emph{Do not use.}
Use $SO(5)$ if you need that level of expressiveness.
\end{result}

\section{Position Signal}
\label{sec:position-signal}

An independent but related mechanism: feeding a position-dependent
scalar $\log(1+t)$ into the gate and lasso projections. This breaks
the shift invariance of the scan, making the network formally dependent
on absolute position.

\begin{result}[Position signal, 2026-04-14]
At scale 8, 10k steps: bpb-neutral. The scan can already break shift
invariance via the lasso's input-dependent rotations; the explicit
position signal adds noise without adding signal. Available as
\texttt{--position-signal} but not recommended on the mainline.
\end{result}

In Session 11 at scale 12 the position signal combined with the lasso
showed a \emph{position miracle}: $\posr$ jumped from 0.001 to 0.275. This
is discussed in detail in Chapter \ref{ch:session11}.

% ------------------------------------------------------------
\chapter{Readouts}
\label{ch:readouts}

The readout $f_\theta: \mathbb{R}^M \to \mathbb{R}^V$ is where most of the
parameters live and where most of the bpb comes from. Chapter
\ref{ch:forensics-omega} will make this concrete: changing the readout
from MLP to routed-experts reliably improves bpb by $\sim 0.02$ and,
more importantly, unlocks the gradient path to $\omegaproj$.

\section{MLP Readout}

A single hidden layer:
\begin{equation}
  f_\theta(h) = W_2 \,\sigma(W_1 h + b_1) + b_2,
  \quad W_1 \in \mathbb{R}^{h \times M},\ W_2 \in \mathbb{R}^{V \times h}
\end{equation}
with $\sigma$ typically GELU, and $h$ an intermediate dimension ($h \approx
256$ at scale 8). The MLP is the default readout (\texttt{--linear-readout-kind
mlp}) and has been the baseline for all 14 Session 11 \texttt{mut-*}
experiments.

\section{Routed Squared-ReLU Experts}

A learned router assigns each token to $k$ of $N$ experts:
\begin{equation}
  f_\theta(h) = \sum_{j=1}^N g_j(h) \cdot E_j(h),
  \quad g = \mathrm{softmax}(W_r h),\ E_j(h) = U_j\,\sigma_{\text{sqrelu}}(V_j h)
\end{equation}
where $\sigma_{\text{sqrelu}}(x) = \mathrm{ReLU}(x)^2$ (the squared ReLU
activation), and $V_j, U_j$ are expert-specific projections.
\texttt{--linear-readout-kind routed\_sqrelu\_experts} with
\texttt{--linear-readout-num-experts 2} is the most thoroughly tested
configuration.

The router softmax breaks the zero-init symmetry that, as we will see in
Chapter \ref{ch:forensics-omega}, blocks $\omegaproj$ gradient in the MLP
readout. Heinrich Session 11 confirms: routed readout has $\|\omegaproj\|_2
\approx 250$ after 50k steps; MLP readout has $\|\omegaproj\|_2 = 0.00000$.

\section{Tied-Embedding Readout}

A Session 8 innovation: the readout projects to embedding space, then
scores via the transposed input embedding.
\begin{equation}
  f_\theta(h) = E^\top \cdot \mathrm{MLP}(h),
  \quad E \in \mathbb{R}^{V \times d}
\end{equation}

Since $E$ is already trained by the input path, it does not need to be
re-learned by the readout. At sp8192 this delivered a $73\%$ parameter
reduction in the readout. On bytes ($V = 256$) the savings are less
dramatic but still meaningful: $5$\,M params instead of $18$\,M for the
same substrate dimension.

\section{Bands}

The \texttt{--readout-bands} flag splits the substrate into contiguous
bands and runs an independent readout per band. Band 1 covers
short-half-life modes (content); band 2 covers medium; band 3 and 4
cover long-context. Session 7 showed that bands differentiate: fast
bands route on context, slow bands route on identity, and band 3 is often
linearly computable from the trained weights.

A common configuration: \texttt{--readout-bands 4 --num-blocks 2}. We use
\texttt{--readout-bands 1} in the Session 11 mut family because the band
mechanism compounds with readout changes in ways we wanted to measure
separately.

% ------------------------------------------------------------
\chapter{Adaptive Substrate Transforms}
\label{ch:adaptive-transforms}

The adaptive substrate introduced in Chapter \ref{ch:causal-bank} admits
six primitives introduced in Session 8 that modify the state update
between the EMA and the readout. All are off by default; all are toggled
by explicit flags.

\section{Write-Side Primitives}

\subsection{Substrate Bank Router}
\texttt{--substrate-bank-router}. Routes tokens to mode banks based on
token type. Adds $\sim 2$k parameters. Effect: small ($< 0.01$ bpb).
Used in combination with other primitives.

\subsection{Overwrite Gate}
\texttt{--overwrite-gate}. Per-mode sigmoid gate that erases stale state
when a new input indicates a topic shift. Adds $\sim 130$k parameters.
Effect: $-0.015$ bpb at s8 in isolation, less when combined with persistent
state.

\section{Read-Side Primitives}

\subsection{Magnitude Normalization}
\texttt{--magnitude-normalize}. $L_2$-normalizes the substrate before
readout, eliminating the ``position counter'' that emerges when long
half-life modes accumulate monotonically. 1 parameter (a learned scale).
Effect: negligible on bpb but improves training stability.

\subsection{Mode Selector}
\texttt{--mode-selector}. Per-token soft attention over modes. Learns
which modes are relevant for a given token. $\sim 130$k parameters.
Effect: small in isolation.

\subsection{Temporal Attention}
\texttt{--temporal-attention}. Cross-attention over substrate snapshots
at fixed intervals. $\sim 50$k parameters. Effect: small; the substrate's
EMA already does most of this work.

\subsection{Hash Memory}
\texttt{--hash-memory} (Session 11 patch). A hash-indexed associative
memory where each embedding is hashed and retrieves a learned vector. O(1)
lookup per token. Prior to commit \texttt{50b5f3d} this module was
\emph{silently bypassed} by the adaptive substrate forward path --- its
gradient was zero for 20k steps of training. The fix plumbed it into both
\texttt{\_linear\_logits} and \texttt{\_forward\_raw}; gradient is now
visible with max $|\mathrm{grad}| \approx 17$.

\section{The Silent Bypass Lesson}

Session 11 discovered that when the \texttt{--adaptive-substrate} flag is
on, the forward path takes an early return in two places ---
\texttt{\_linear\_logits} and \texttt{\_forward\_raw} --- that skip most
of the write-side primitives. This was silent: no error, no warning,
just zero gradient to the skipped modules. The hash-memory plugin was
trained for 20k steps while receiving no gradient, and every experiment
on that branch was bpb-neutral despite showing the hash-memory flag in
the command line.

The lesson, formalized as a project-level feedback rule:

\begin{quote}
\emph{Always plumb new modules into BOTH adaptive and non-adaptive forward
branches; gradient-test before training.} When adding a module that
interacts with the substrate, write a gradient-sign fixture: set the
module's output to have a clear effect on the loss, backprop once, assert
the module's weights have non-zero gradient.
\end{quote}

Three gradient-sign fixtures now live in
\texttt{decepticons/tests/test\_adaptive\_substrate\_plumbing.py} (commit
\texttt{3eaf0aa}):
\begin{enumerate}
  \item Hash memory receives gradient when \texttt{adaptive\_substrate=True}.
  \item Local path receives gradient under adaptive substrate.
  \item Local path is properly disabled when \texttt{--local-scale-override 0.0}.
\end{enumerate}

These run on every decepticons commit and would have caught the silent
bypass in Session 11 if they had existed in Session 10.

% ============================================================
\part{The Chronicle: Sessions 6 through 11}
% ============================================================

\chapter{Session 6: The Causal Bank}
\label{ch:session6}

Session 6 (\texttt{paper/session6\_chronohorn.tex}) introduced the causal-bank
family and the measurement methodology in Chapter \ref{ch:methodology}.
Key results:
\begin{itemize}
  \item Frozen substrate at scale 8, 10k steps, sp8192: 1.92 bpb with
    MLP readout. This established the baseline.
  \item Half-life distribution shown to be load-bearing: $\tau_{\max}$
    from 128 to 1024 moves bpb by $\sim 0.04$.
  \item Causality test introduced; one version of the scan was caught
    leaking from $t+1$ into $t$ via a Triton bug. Fixed before merge.
\end{itemize}

\chapter{Session 7: The Bottleneck}
\label{ch:session7}

Session 7 (\texttt{paper/session7\_bottleneck.tex} and the companion
\texttt{session7\_five\_papers.tex}) identified that the mainline causal-bank
was capacity-limited not by parameters but by \emph{content coupling}: the
linear projection $W_{\text{in}}$ in Equation \ref{eq:substrate} was the
bottleneck. Three things came out of Session 7:
\begin{enumerate}
  \item The \texttt{--input-proj-scheme} flag with the \texttt{random}
    option: at init, fill $W_{\text{in}}$ with i.i.d. Gaussian rather
    than Xavier. This breaks the symmetry that a learned $W_{\text{in}}$
    has trouble escaping during the first epoch.
  \item The adaptive substrate path (Chapter \ref{ch:causal-bank}),
    making $W_{\text{in}}$ effectively input-dependent via $\omegaproj$.
  \item The bands readout (Chapter \ref{ch:readouts}), which splits the
    readout across different half-life ranges.
\end{enumerate}

Session 7 also introduced the split into decepticons and chronohorn as
separate codebases. By the end of Session 7, the best model on sp8192 was
1.78 bpb with bands + adaptive substrate at scale 10.

\chapter{Session 8: Twelve Hours, Forty Experiments}
\label{ch:session8}

Session 8 (\texttt{paper/session8\_chronohorn.tex}) was a stress test of
the fleet. Forty experiments were dispatched in twelve hours; every
mechanism in Chapter \ref{ch:adaptive-transforms} was introduced,
independently ablated, and logged to the database. The central finding:
most of the read-side primitives delivered within noise of each other,
which is an informative result --- the bottleneck was no longer on the
read side.

Tied-embedding readout was introduced in Session 8 and validated at sp8192:
$73\%$ parameter reduction in the readout for no bpb loss.

Session 8 also introduced the ``firewall test'' that enforces the
\texttt{chronohorn \not\to decepticons} dependency direction. At the
start of Session 8 there were seven cross-imports; all were removed.

\chapter{Session 9: The Lasso and the Lie Algebra}
\label{ch:session9}

Session 9 (\texttt{paper/session9\_chronohorn.tex}; introduced by the
commit message ``the lasso, the Lie algebra, and 0.001'') was the rotation
session. All five rotation primitives in Chapter \ref{ch:rotations} were
introduced, implemented, and benchmarked:
\begin{center}
\begin{tabular}{lrr}
\toprule
Rotation & $\Delta$ bpb vs baseline & tok/s \\
\midrule
None (baseline)       &  0.000 & 316k \\
Complex (SO(2))       & $-0.008$ & 312k \\
Lasso ($2\times 2$)   & $\mathbf{-0.109}$ & 316k \\
Quaternion (SO(3))    & $-0.107$ & 294k \\
SO(5)                 & $-0.105$ & 117k \\
Quintic (GL(5))       & NaN & --- \\
\bottomrule
\end{tabular}
\end{center}

Key findings:
\begin{itemize}
  \item Position $R^2 = 0.001$ across all rotations. Order is not the
    bottleneck at this bpb range.
  \item Content $R^2$ \emph{doubled} (0.101 to 0.209) with the lasso.
    The rotation is doing its work through richer content coupling,
    not order encoding.
  \item The lasso scan can be fused via Triton. The Hillis--Steele
    parallel prefix scan with \texttt{torch.roll} and \texttt{torch.where}
    is dynamo-compatible.
\end{itemize}

Session 9 also initiated the byte-level pivot. Until Session 9 the
mainline was sp1024; Session 10 would finish the pivot.

\chapter{Session 10: Structure Solved}
\label{ch:session10}

Session 10 (\texttt{paper/session10\_byte\_frontier.tex} and the companion
\texttt{session10\_structure\_solved.tex}) moved the mainline to byte level
and pushed the frontier to 1.785 bpb with the \emph{learnable-s8}
configuration: routed\_sqrelu\_experts readout, adaptive substrate with
\texttt{--hrr-omega-init}, \texttt{--adaptive-shared-proj}, scale 8, 50k
steps.

Two mechanisms dominated:
\begin{itemize}
  \item \emph{Learnable $\omegaproj$.} The projection that maps input
    embeddings into mode-frequency modulations was allowed to train
    rather than be frozen at the Fourier basis.
  \item \emph{Shared-proj speed flag.} One trunk matmul
    (embed\_dim $\to$ modes) plus five cheap head matmuls, replacing
    five large independent projections. Brought throughput from 26k
    tok/s (Session 9) to 300k+ tok/s.
\end{itemize}

Session 10 also discovered the \emph{scan is dead} result: $\cka = 1.000$
across all checkpoints, half-life effective below 6, gradient vanishes
through the scan. The EMA substrate was delivering no signal beyond what
the local path provided. This was confirmed by ablating the scan entirely
and finding bpb change of $< 0.01$.

The implication: the substrate's value is in $\omegaproj$, not the scan.
The scan is infrastructure; the modulation is the substance.

A second finding: at scale 12 with the lasso and position signal, $\posr$
jumped from 0.001 (Session 9) to 0.275. \emph{Something} was learning
about position. Session 11 would investigate this.

\chapter{Session 11: Mechanism Decomposition}
\label{ch:session11}

Session 11 is the current session as of this writing. Fourteen \texttt{mut-*}
experiments, all at s8 with MLP readout, all within 0.015 bpb of each
other (band 2.10--2.12). This clustering is diagnostic: \emph{the MLP
readout has a ceiling at ~2.10 bpb}, regardless of which write-side or
read-side primitive is added.

\section{The Heinrich Readout}

Heinrich's MRI of five critical models:
\begin{center}
\begin{tabular}{lrrrrr}
\toprule
Model & params & bpb & $\effdim$ & $\conr$ & $\|\omegaproj\|_2$ \\
\midrule
baseline s8 (MLP, frozen)    & 17.8\,M & 2.075 & 150 & 0.342 & $0$ \\
mut-control-mlp              & 17.8\,M & 2.051 & 150 & 0.359 & $0$ \\
mut-persistent (MLP)         & 17.8\,M & 2.054 & 148 & 0.345 & $0$ \\
mut-routed (routed, frozen)  & 17.8\,M & 2.049 & 161 & 0.332 & $0$ \\
mut-mlp-omega10 (MLP, 10x LR)& 17.8\,M & 2.057 & 151 & 0.359 & $0$ \\
\textbf{learnable-s8} (routed, $\omegaproj$) & 17.8\,M & \textbf{1.785} & \textbf{314} & 0.350 & \textbf{250} \\
\bottomrule
\end{tabular}
\end{center}

The contrast between \texttt{mut-routed} (routed readout, $\omegaproj$ at
zero) and \texttt{learnable-s8} (routed readout, $\omegaproj$ trained to
$\|\omegaproj\|_2 = 250$) is the single most important observation in
this session. The readout and $\omegaproj$ together double $\effdim$ from
150 to 314 and cut bpb by 0.26.

\section{The Silent Structural Block}

Heinrich's answer to the first of six questions in Chapter
\ref{ch:forensics-omega} revealed that \texttt{mut-mlp-omega10} --- which
set a $10\times$ learning rate multiplier on $\omegaproj$ --- had
\emph{literally zero gradient} at every probe: max $|\omegaproj| =
0.00000$ at steps 5k, 10k, 15k, and 20k. The $10\times$ LR was applied
to updates of magnitude $0$; $10 \times 0 = 0$.

Why: $\omegaproj$ is initialized to zero (the bias has the HRR Fourier
basis, the weight is zero). For the MLP readout, the gradient chain back
to $\omegaproj$ must pass through a single-hidden-layer $\sigma(W_1 h +
b_1)$. At init, the relevant Jacobian entries are zero, and the chain
multiplies to zero. The routed readout's softmax breaks this symmetry
because $\partial g / \partial h$ is non-zero at init (the softmax is
non-constant).

This is a \emph{structural} gradient block, not an optimization one. It
cannot be escaped by changing LR; the only escape is to change the
readout or to noise-initialize $\omegaproj$.

\section{The Orthogonal Mechanism}

Session 11 also introduced \emph{gated-delta-cr}, a substrate that
combines gated-delta dynamics with complex rotation. At s8 it achieves
2.074 bpb at 5.7\,M parameters --- less than a third of the MLP-family
parameters for better bpb. Heinrich's MRI:
\begin{itemize}
  \item $\effdim = 24$ (out of 2048). Aggressive compression.
  \item 94\% classification accuracy on the current byte from the top
    24 PCs. The substrate is a \emph{byte-identity encoder}.
  \item $\posr < 0.001$. Zero position information in the substrate.
  \item $\cka$ with baseline adaptive substrate $= 0.21$. Orthogonal
    subspace.
\end{itemize}

This is not a variant of the mainline; it is a different architecture.
Context in gated-delta-cr lives in the rotation phases between states,
not in an accumulated substrate. Scaling gated-delta-cr to s12 delivered
\emph{worse} bpb (2.126, $+0.05$ from s8) --- confirming that its
mechanism doesn't absorb scale the way the adaptive substrate does.

\section{The Current Frontier}

At the end of Session 11 ablations (still running as of this book):
\begin{itemize}
  \item \textbf{mut-lasso-pos-b2-s8-50k}: 2.048 bpb. Extending training
    time from 20k to 50k at s8 (not s12) cleaned up the position signal
    and dropped bpb by 0.039. Scale was not required.
  \item \textbf{mut-routed-learnable-persistent-s12-50k}: running at
    step 25.6k/50k with probe bpb 1.939; projecting 1.80--1.82 at
    50k per the additive-law prediction.
  \item \textbf{learnable-s8}: 1.785 bpb. Still the frontier for byte-level
    mainline architectures.
\end{itemize}

% ============================================================
\part{Forensics: The Heinrich Toolkit}
% ============================================================

\chapter{The Instruments}
\label{ch:forensics}

Heinrich is a set of tools for opening a trained model and asking
specific questions. All tools take a checkpoint and some test data and
return a summary. The five most-used:

\begin{description}
  \item[$\effdim$] Effective dimensionality. Given a set of activations
    $H \in \mathbb{R}^{n \times d}$, let $\lambda_i$ be the eigenvalues
    of $H^\top H$. Then
    \begin{equation}
      \effdim(H) = \frac{(\sum_i \lambda_i)^2}{\sum_i \lambda_i^2}.
    \end{equation}
    Interpretation: the number of directions doing work. For a 2048-mode
    substrate, $\effdim \in [1, 2048]$; typical trained models are in
    the 100--300 range.
  \item[$\conr$] Content $R^2$. Linear regression of the substrate on
    the current-byte one-hot; the $R^2$ of that regression. Measures
    how much of the substrate can be linearly read out as ``what byte
    am I?''.
  \item[$\posr$] Position $R^2$. Linear regression of the substrate on
    one-hot position, then the $R^2$.
  \item[$\cka$] Centered Kernel Alignment between two sets of activations.
    Scalar in $[0, 1]$ measuring how similar the two subspaces are.
    $\cka = 1$ means same subspace; $\cka = 0$ means orthogonal.
  \item[$\sigma(\omegaproj)$] Singular value spectrum of the
    $\omegaproj$ matrix. Flat means full-rank; a spectrum concentrated
    at the top means low-rank.
\end{description}

\section{Why Linear Regression?}

The $\conr$ and $\posr$ metrics use linear regression, not an MLP. This
is intentional. A linear regression measures how much of the target
variable is \emph{linearly decodable} from the substrate. That is
precisely the information that a linear readout can use. If a nonlinear
regression has higher $R^2$, the readout is not seeing that information
--- it's hidden behind nonlinearity that the readout doesn't have.

Chapter \ref{ch:session11} contains the most striking consequence of this
choice: the gated-delta-cr substrate has $\posr < 0.001$ but gets 2.07
bpb anyway. Its context information is in the \emph{dynamics} of the
rotation phases, not in the linearly-decodable substrate state. A linear
probe is blind to this; a full forward evaluation through the model is
not.

% ------------------------------------------------------------
\chapter{The $\omegaproj$ Dead Gradient}
\label{ch:forensics-omega}

The longest-lived bug in decepticons' modeling code was the MLP-readout
gradient block described in Section \ref{ch:session11}. This chapter
presents it formally.

\section{The Symmetry}

Consider the adaptive substrate path at initialization:
\begin{align}
  \omega_t &= \omegaproj x_t,\quad \omegaproj \in \mathbb{R}^{M \times d},\ \omegaproj = 0, \\
  h_t &= f_{\text{HRR}}(W_{\text{in}} x_t, \omega_t), \\
  h_t &\stackrel{\text{init}}{=} f_{\text{HRR}}(W_{\text{in}} x_t, 0) = W_{\text{in}} x_t \odot v_0
\end{align}
where $v_0$ is the HRR Fourier basis in $\omegaproj$'s bias. At init,
$\omegaproj$'s weight is zero, so the modulation is entirely from the
frozen Fourier basis.

Now consider the gradient. Let $\ell = \mathrm{CE}(y, f_\theta(h))$ be
the cross-entropy loss. Then
\begin{equation}
  \frac{\partial \ell}{\partial \omegaproj} = \frac{\partial \ell}{\partial \omega}\cdot x^\top,
  \quad
  \frac{\partial \ell}{\partial \omega} = \frac{\partial \ell}{\partial h}\cdot\frac{\partial h}{\partial \omega}.
\end{equation}

The first factor, $\partial \ell / \partial h$, flows back through the
readout. For the MLP readout at init:
\begin{equation}
  f_\theta(h) = W_2\,\sigma(W_1 h + b_1) + b_2,
  \quad \frac{\partial f}{\partial h} = W_2\,\sigma'(\cdot)\,W_1.
\end{equation}

With standard initialization, $W_1 h$ is symmetric in expectation, and
since $\sigma$ (GELU, ReLU) has a zero-activation neighborhood, many
Jacobian entries are zero. The cascade of Jacobians yields a
$\partial \ell / \partial h$ that --- when combined with
$\partial h / \partial \omega$, which itself is zero at $\omegaproj = 0$
--- produces a gradient of zero.

\begin{result}[MLP structural block]
For the MLP readout with $\omegaproj = 0$ at initialization and GELU
activation, $\partial \ell / \partial \omegaproj = 0$ to machine
precision, and this state is an attractor: any finite-precision
perturbation does not escape.
\end{result}

Heinrich's measurement: at every checkpoint of \texttt{mut-mlp-omega10}
(5k, 10k, 15k, 20k), $|\omegaproj|_\infty = 0.00000$ bit-identically.
The $10\times$ LR multiplier, the $0.1$ weight decay, the cosine
schedule --- none of it matters. Zero times any scalar is zero.

\section{The Routed Escape}

The routed-experts readout has a softmax router:
\begin{equation}
  g(h) = \mathrm{softmax}(W_r h).
\end{equation}
At init, $W_r h$ may be small but the softmax is non-constant: its output
is not uniformly distributed across experts (random weight init breaks
this symmetry). So $\partial g / \partial h$ is non-zero, and the gradient
chain does reach $\omegaproj$.

Heinrich's measurement: routed-readout $\omegaproj$ at 50k steps has
$\|\omegaproj\|_2 = 250$. The singular value spectrum is heavily
low-rank (see below), but the gradient did flow.

\section{The Rank of a Trained $\omegaproj$}

For \texttt{byte-hrr-learnable-s8-50k}:
\begin{itemize}
  \item Shape: $(2048, 256)$.
  \item $\|\omegaproj\|_2 = 249.6$.
  \item Top-10 singular values: 73.1, 70.7, 66.4, 64.4, 62.6, 60.3,
    56.8, 53.8, 52.5, 50.3.
  \item Tail-10 singular values: $\leq 0.5$.
  \item Rank to 50\% energy: 8 / 256.
  \item Rank to 90\% energy: 22 / 256.
  \item Rank to 99\% energy: 63 / 256.
  \item Effective rank (entropy): 28.66.
\end{itemize}

The matrix is genuinely low-rank. Of 256 input directions, only $\sim
29$ participate meaningfully in driving the 2048 mode-frequency
modulations. This suggests a design-space simplification: parameterize
$\omegaproj = U V^\top$ with $U \in \mathbb{R}^{256 \times 32}$, $V \in
\mathbb{R}^{2048 \times 32}$. Parameter count drops from 524k to 73k,
a $7\times$ savings. This is an open experiment at the time of writing.

\section{The Orthogonal Diagnosis}

The ``orthogonal path'' of gated-delta-cr is visible in the Heinrich
readout:
\begin{center}
\begin{tabular}{lrrr}
\toprule
Model & $\effdim$ & $\conr$ & $\posr$ \\
\midrule
gated-delta-cr-s8 & 24 & 0.104 & $< 0.001$ \\
learnable-s8      & 314 & 0.350 & 0.003 \\
baseline-s8       & 150 & 0.342 & 0.001 \\
\bottomrule
\end{tabular}
\end{center}

gated-delta-cr has one-fifth the effective dimensionality of learnable-s8
but delivers 2.07 vs 1.785 bpb. It is \emph{not} using the substrate the
way learnable-s8 is; the context information lives in the rotation
dynamics, not in the linearly-readable substrate state. That this delivers
competitive bpb at a fraction of the parameters is the most interesting
positive result of Session 11.

% ============================================================
\part{Infrastructure}
% ============================================================

\chapter{The Database}
\label{ch:database}

Chronohorn's database is a single SQLite file at \texttt{out/chronohorn.db}.
It is the source of truth for every result, probe, job, and event in the
project. Raw JSON archives of results are preserved on Sharts but are
treated as a backup: all dashboards, MCP tools, and the frontier query
read from the DB.

\section{Schema}

The database has 11 tables:

\begin{itemize}
  \item \texttt{results}: one row per completed training run.
    Columns: \texttt{name}, \texttt{config\_id}, \texttt{bpb},
    \texttt{train\_bpb}, \texttt{tok\_s}, \texttt{total\_tflops},
    \texttt{wall\_sec}, \texttt{slope}, \texttt{illegal},
    \texttt{eval\_batches}, \texttt{accelerator\_arch}, and a few others.
  \item \texttt{configs}: one row per unique configuration (hash-keyed).
    Columns include substrate\_mode, readout, scale, state\_dim, etc.,
    plus a full JSON blob.
  \item \texttt{probes}: one row per probe, keyed by (\texttt{name},
    \texttt{step}). Columns: \texttt{bpb}, \texttt{loss},
    \texttt{elapsed\_sec}, \texttt{train\_elapsed\_sec},
    \texttt{eval\_batches}.
  \item \texttt{jobs}: one row per dispatched job. Columns:
    \texttt{state}, \texttt{launched\_at}, \texttt{completed\_at},
    \texttt{host}, \texttt{remote\_run}, \texttt{runtime\_job\_name},
    \texttt{failure\_reason}, checkpoint path, etc.
  \item \texttt{forecasts}: per-result fit parameters from a saturation
    analysis (e.g. power-law asymptote).
  \item \texttt{predictions}: model-predicted bpb for jobs not yet run.
  \item \texttt{events}: time-series log of drain ticks, launches,
    completions, errors, catchup fires, reconciliations.
  \item \texttt{fleet}: latest probe of each host (GPU util, memory,
    online, containers).
  \item \texttt{journal}: free-form experiment journal entries.
  \item \texttt{schema\_version}: single-row table holding the current
    schema version; migrations run on open.
  \item \texttt{sqlite\_sequence}: SQLite internal.
\end{itemize}

\section{Single-Writer Discipline}

SQLite supports one writer at a time. Chronohorn enforces this via a
dedicated writer thread: all mutations go through a queue, and
\texttt{\_write} / \texttt{\_write\_many} wait for the queue to drain
when \texttt{wait=True}.

Reads go through a separate connection in autocommit WAL mode, so readers
never block writers and vice versa. The read lock is a module-level
\texttt{threading.Lock} that serializes reads against the writer thread's
periodic schema-check queries.

\section{Schema Versioning}

At database open, the code compares the current schema version to the
target. Migrations run sequentially (v1 $\to$ v2 $\to \ldots$). The
current version is 15. Recent migrations:
\begin{itemize}
  \item v14: added \texttt{probes.train\_elapsed\_sec} to distinguish
    training time from wall time across probe breaks.
  \item v15: added fleet GPU telemetry columns
    (\texttt{gpu\_util\_pct}, \texttt{gpu\_mem\_used\_mb},
    \texttt{gpu\_mem\_total\_mb}).
\end{itemize}

\chapter{The Runtime Daemon}
\label{ch:runtime}

\texttt{chronohorn runtime} is the unified daemon. One process runs
three long-running threads plus an HTTP server:
\begin{enumerate}
  \item \texttt{fleet\_probe}: 30-second cadence. SSH to each host,
    read \texttt{nvidia-smi} and \texttt{kubectl get pods}, write
    \texttt{fleet} row.
  \item \texttt{drain}: 60-second cadence. Materialize manifest jobs
    into \texttt{jobs} table, dispatch pending to available GPUs,
    pull completed results, reconcile DB state against k8s state,
    clean up ghosts.
  \item \texttt{catchup}: 120-second cadence. Rescue completed or
    failed jobs that are missing their Sharts exports. This used to
    run inside \texttt{drain} but blocked the loop for 10+ minutes
    during large catch-up batches; it was split out at the end of
    Session 11.
  \item \texttt{http}: synchronous handler for \texttt{/api/status},
    \texttt{/api/action}, \texttt{/api/events}, and the dashboard HTML.
\end{enumerate}

Flags:
\begin{itemize}
  \item \texttt{--manifest}: manifest JSONL paths to drain. Repeatable.
    Without any, drain processes all active DB jobs.
  \item \texttt{--no-dispatch}: monitor and pull only. Pending jobs
    stay pending.
  \item \texttt{--auto-deepen}: auto-generate deepening runs when slope
    is alive.
  \item \texttt{--poll N}: drain poll interval in seconds.
  \item \texttt{--no-browser}: skip Chrome app launch.
\end{itemize}

\chapter{The Fleet Drain Loop}
\label{ch:drain}

The drain loop is where the complexity of distributed experimentation
lives. Given a manifest, a database, and a fleet, the drain must answer
three questions on every tick:
\begin{enumerate}
  \item Which jobs are pending and ready to launch?
  \item Which jobs are running and producing probes?
  \item Which jobs are completed and need their results pulled back?
\end{enumerate}

The drain answers these by calling \texttt{probe\_fleet\_state} for a
live view of the cluster and \texttt{db.active\_jobs()} for a database
view. It partitions the union into pending / running / completed / stale,
launches pending into free slots, pulls completed into the DB, and
reconciles DB state against fleet state.

\section{Partition}

\texttt{partition\_running\_jobs(jobs, fleet\_state)} classifies each
job:
\begin{itemize}
  \item Running: k8s pod is active and GPU is busy.
  \item Completed: k8s Job completed successfully; result JSON should
    be available on the host path.
  \item Stale: k8s Job is running but GPU is idle past a threshold
    (training is hung).
  \item Pending: otherwise.
\end{itemize}

\section{Ghost Cleanup}

Jobs that are \texttt{dispatched} or \texttt{running} in the DB for
more than 6 hours, and that the fleet cannot see (pod is gone), are
\emph{ghosts}. Until Session 11, ghost cleanup flipped them to
\texttt{state='failed'} with \texttt{failure\_reason='ghost\_cleanup'}.

This produced a critical false-positive. Three pilots
(\texttt{pilot-A-triton-default-1500},
\texttt{pilot-C-batch64-1500},
\texttt{pilot-D-lowrank-1500}) completed successfully but had their
k8s Job records TTL-cleaned after 5 minutes. The next drain tick saw
them as ghosts and marked them failed. The result JSONs sat on the
host filesystem for 12+ hours before the catch-up logic was widened
to include failed-state jobs.

The fix (Session 11, commit \texttt{4acb6ce}): before marking a ghost
as failed, check whether the \texttt{results} table has a row for that
name. If yes, the job actually completed; flip to \texttt{state='completed'}
and clear \texttt{failure\_reason}. Only ghosts without results get
marked failed.

\section{The Catch-Up Path}

The catch-up path (commits \texttt{7392443} and \texttt{4acb6ce}) scans
for jobs whose \texttt{state} is \texttt{completed} or \texttt{failed}
and whose Sharts export JSON is missing. For each, it invokes
\texttt{pull\_remote\_result}, which is idempotent: skips the SCP if
the local JSON already exists, but runs the Sharts export unconditionally.

The catch-up widening to state='failed' was necessary because
ghost-cleanup's false positives never got their results pulled until
catch-up explicitly targeted them. Before the widening, the only way to
recover a false-ghosted result was manual SSH and SCP.

The catch-up runs in a dedicated daemon thread now (Session 11, same
commit), on a 120-second cadence, so that large SCP batches don't block
drain responsiveness. A batch of 25 jobs $\times$ 5 checkpoints $\times$
70 MB each is $\sim 8.75$ GB of SCP traffic, which can take 10+ minutes
over a LAN. Blocking the drain loop for that long meant no new launches,
no reconciliation, and fleet probes piling up. Separate thread fixes all
three.

\chapter{The Auto-Pull Pipeline}
\label{ch:autopull}

\section{Design}

\begin{center}
\begin{tikzpicture}[node distance=0.6cm and 1.8cm,
  box/.style={rectangle,draw,rounded corners=2pt,minimum width=3cm,minimum height=0.8cm,align=center,font=\footnotesize},
  arrow/.style={->,>=Stealth,thick}]
\node[box] (k8s) {k8s Job on slop-XX};
\node[box, below=of k8s] (pod) {pod writes /run/results/\texttt{<name>.json}};
\node[box, below=of pod] (host) {hostPath /tmp/chronohorn-runs/\texttt{<name>}/results/};
\node[box, below=of host] (drain) {drain.pull\_remote\_result};
\node[box, below=of drain] (db) {results table};
\node[box, right=of drain, xshift=1cm] (sharts) {\_export\_checkpoint \\to Sharts/session11/};
\draw[arrow] (k8s) -- (pod);
\draw[arrow] (pod) -- (host);
\draw[arrow] (host) -- (drain);
\draw[arrow] (drain) -- (db);
\draw[arrow] (drain) -- (sharts);
\end{tikzpicture}
\end{center}

\section{Bugs and Their Fixes}

Over the course of Session 11 the auto-pull pipeline had six discrete
bugs, each discovered when it produced a stranded or silently-failed
result. A partial chronicle:

\begin{enumerate}
  \item \emph{Daemon auto-respawn.} A running
    \texttt{chronohorn runtime --manifest} daemon will re-dispatch jobs
    that are killed via the fleet or manually. Fix: check for the
    daemon process before killing jobs, or use \texttt{--no-dispatch}
    mode.
  \item \emph{Checkpoint pull only the end-of-training file.} Periodic
    checkpoints saved by \texttt{--save-checkpoint-every N} were
    silently dropped. Fix: enumerate and SCP all
    \texttt{<name>\_step*.checkpoint.pt} files alongside the final.
  \item \emph{MCP \texttt{chronohorn\_fleet\_pull} missed per-job
    launches.} Legacy path scraped
    \texttt{/data/chronohorn/out/results/*.json}; per-job launches
    write to \texttt{/tmp/chronohorn-runs/<job>/results/}. Fix: MCP
    tool now runs both legacy and per-job paths.
  \item \emph{Catch-up only considered state='completed'.} Pilots
    ghost-cleaned to state='failed' never got their results pulled.
    Fix: widen catch-up to state in (completed, failed) with
    Sharts-exists skip gate.
  \item \emph{NaN bpb crashed drain.} A result JSON with
    \texttt{test\_bpb = nan} passed the ``non-null'' check,
    slipped through the ``$\leq 0$'' positivity check (NaN compares
    false both ways), and hit SQLite as NULL $\to$ NOT NULL
    constraint. Fix: explicit \texttt{math.isnan(bpb)} guard in
    \texttt{record\_result}; widen ingestion except to catch
    IntegrityError.
  \item \emph{Catch-up blocked the drain tick.} A batch of 25 jobs
    with 5 checkpoints each made the tick take 10+ minutes. Fix:
    split catch-up into its own daemon thread at 120s cadence.
\end{enumerate}

All six fixes are covered by regression tests in
\texttt{tests/test\_drain.py} and \texttt{tests/test\_result\_pullback.py}.

\chapter{The MCP Surface}
\label{ch:mcp}

Chronohorn exposes 55 tools over MCP (Model Context Protocol). The
surface is the primary way an AI collaborator interacts with the project:
not by calling Python functions directly, but by invoking named tools
that return JSON.

\section{Categories}

\begin{itemize}
  \item \emph{Manifests and results.}
    \texttt{chronohorn\_manifests} ingests a manifest;
    \texttt{chronohorn\_results} ingests result JSONs;
    \texttt{chronohorn\_list\_manifests} enumerates known manifests.
  \item \emph{Fleet.}
    \texttt{chronohorn\_fleet\_status} probes hosts;
    \texttt{chronohorn\_fleet\_pull} syncs results;
    \texttt{chronohorn\_fleet\_dispatch} launches pending jobs;
    \texttt{chronohorn\_fleet\_sync} does pull + probe + frontier in one.
  \item \emph{Query.}
    \texttt{chronohorn\_query} runs a raw SELECT;
    \texttt{chronohorn\_frontier} returns the admissible frontier;
    \texttt{chronohorn\_compare} compares named runs.
  \item \emph{Analysis.}
    \texttt{chronohorn\_learning\_curve},
    \texttt{chronohorn\_saturation},
    \texttt{chronohorn\_predict},
    \texttt{chronohorn\_forecast}.
  \item \emph{Control.}
    \texttt{chronohorn\_stop\_run},
    \texttt{chronohorn\_k8s\_delete\_job},
    \texttt{chronohorn\_reset} (poisoned-DB purge).
  \item \emph{Meta.}
    \texttt{chronohorn\_changelog} summarizes DB history;
    \texttt{chronohorn\_journal\_write} writes session journal.
\end{itemize}

The handler routes \texttt{call\_tool(name, args)} to a dispatch table
keyed on tool name. Adding a tool is a one-liner in the dispatch dict
plus a method; no registration ritual.

\chapter{The Dashboard}
\label{ch:dashboard}

\texttt{chronohorn observe serve} (now folded into the runtime daemon
as its HTTP server) serves a single-page dashboard at the daemon's port.
The dashboard has eight panes:
\begin{itemize}
  \item \textbf{jobs}: running / dispatched / pending / completed /
    failed rows with state badge, progress, bpb.
  \item \textbf{frontier}: top 30 legal results by bpb.
  \item \textbf{curves}: probe bpb over training step, one curve per
    prefix, log-scaled step axis.
  \item \textbf{eff}: efficiency-ranked results by $\mu\bpb/\text{TFLOP}$
    (marginal bpb improvement per total TFLOP).
  \item \textbf{config}: substrate, readout, bands, scale, params.
  \item \textbf{fleet}: host status, GPU util, active containers.
  \item \textbf{manifests}: known manifests with job counts.
  \item \textbf{events}: tail of the events table.
\end{itemize}

The backend is a single \texttt{\_build\_api\_data} function that
executes a small fixed set of DB queries (no N+1 lookups) and returns
a JSON blob. The dashboard receives updates via Server-Sent Events on
\texttt{/api/events}; the runtime pushes a new blob whenever a drain
tick launches, pulls, or ingests probes.

\section{Known Bugs and Their Fixes}

Session 11 fixed five rendering bugs (commit \texttt{8c28a35}):
\begin{enumerate}
  \item Curves pane rendered 700+ historical curves simultaneously.
    Fix: filter to curves from runs with activity in the last 48 hours.
  \item Y-axis labels duplicated (``8.2 8.2 8.4 8.4 \ldots'').
    Fix: dedupe via Set on formatted labels.
  \item $\mu\bpb/\text{TFLOP}$ column showed ``0.0'' for every row
    because values range $10^{-5}$ to $10^{-3}$ and the display
    used \texttt{.toFixed(1)}. Fix: display in micro-bpb units
    (multiply by $10^6$, show integer).
  \item Readout abbreviation ``exp'' was opaque without a legend.
    Fix: render as ``routed-exp''.
  \item Three ghost-pending rows from \texttt{frontier\_hrr} persisted
    on every render because the manifest-active window was 7 days and
    the manifest's last launch was 4 days ago. Fix: tighten
    pending-active window to 3 days; keep completed/failed history at
    7 days.
\end{enumerate}

% ============================================================
\part{The Speed Stack and Paths to 1.0}
% ============================================================

\chapter{Speed: Toward 1 Million Tokens per Second}
\label{ch:speed}

The project target is $10^6$ tok/s on A4000. Session 9 baseline was
26k tok/s on the adaptive substrate. Session 11 has brought it to 324k
tok/s, $12.5\times$ the baseline, still $3\times$ short of the target.

\section{The Tok/s Budget}

At 1M tok/s, $B = 32$, $L = 1024$, the iteration time is 32 ms. For a
forward pass with $\sim 20$ MFLOPs/byte, the compute budget is
$(1e6 \cdot 20e6)\,\text{FLOPs/s} = 20$ TFLOPs, which is the A4000's
FP16 peak. In practice we can hit 60\% of peak, so the real budget is
$\sim 12$ TFLOPs achievable. At 324k tok/s we are achieving $\sim 11$
TFLOPs --- \emph{we are already compute-bound}. Further speedup must
come from doing less work per byte, not from running the same work
faster.

\section{The Moves That Worked}

\begin{itemize}
  \item \emph{Triton fused scan.} The scan was originally a Python loop
    over sequence positions; at $L = 1024$ it was 30\% of forward time.
    Fused into a single Triton kernel, the scan dropped to $<$ 3\%.
    $21\times$ speedup on the scan, $1.8\times$ on end-to-end forward.
  \item \emph{torch.compile default mode.} Session 10 experimented with
    \texttt{reduce-overhead} for CUDA graphs; it was faster on
    synthetic benchmarks but leaked memory on long runs. Default mode
    is stable and $1.5\times$ faster than eager.
  \item \emph{Batch auto-scale to 60\% VRAM.} Chronohorn's
    \texttt{apply\_variant} function calculates the maximum batch
    size that fits in 60\% of VRAM and uses it. At scale 8 on A4000
    this is typically 32.
  \item \emph{Adaptive shared projection.} One trunk matmul
    (embed\_dim $\to$ modes) plus five cheap head matmuls, replacing
    five large independent projections. $3\times$ speedup on the
    adaptive substrate's modulation path.
\end{itemize}

\section{The Moves That Didn't Work}

\begin{itemize}
  \item \emph{reduce-overhead CUDA graphs.} Per-position telemetry writes
    from inside the forward pass corrupted the graph. Skipping the
    telemetry under reduce-overhead fixed the crash but the graph was
    only $\sim 5\%$ faster than default compile, not worth the
    fragility.
  \item \emph{Triton on Turing (slop-home).} Inductor's Triton autotune
    triggers a device-side assert on sm\_75. No known fix; we skip
    slop-home for byte runs.
\end{itemize}

\section{What's Left}

\begin{itemize}
  \item \emph{Fused CUDA kernel for the whole adaptive-substrate path.}
    Currently the path is five Triton calls plus some framework glue.
    A single fused kernel would remove glue overhead; projected
    $\sim 1.5\times$ speedup.
  \item \emph{Low-rank $\omegaproj$.} Chapter \ref{ch:forensics-omega}'s
    rank-29 finding suggests a rank-32 factorization of $\omegaproj$
    is free on bpb and saves $7\times$ parameters. Parameters become
    bandwidth on the A4000; $7\times$ fewer parameters in the
    $\omegaproj$ matrix saves substantial HBM traffic.
  \item \emph{RT core inference.} The A4000's ray-tracing cores are
    idle; a representation of the scan as triangle intersection (the
    rotation as a triangle, the readout as a BVH traversal) would
    give $\sim 37$ TFLOPs unused throughput to the right architecture.
    Open research.
\end{itemize}

\chapter{Paths to 1.0 bpb}
\label{ch:paths}

Current best is 1.785 bpb (learnable-s8 at 50k steps). The gap to 1.0
is 0.78 bpb. Heinrich's Session 11 analysis suggests the following
paths, ranked by expected return:

\begin{enumerate}
  \item \textbf{Combine learnable $\omegaproj$ with persistent state at
    scale 12.} Additive-law prediction: 1.80--1.82 bpb. Pilot
    \texttt{mut-routed-learnable-persistent-s12-50k} is running at the
    time of this writing, probe 1.939 at 25.6k/50k.
  \item \textbf{Scale 14.} All successful mechanisms (lasso,
    learnable $\omegaproj$, persistent state) compound with scale.
    Extrapolation to s14: 1.65 bpb. Cost: $\sim 8\times$ compute vs s8.
  \item \textbf{Low-rank $\omegaproj$.} Free on bpb (Chapter
    \ref{ch:forensics-omega}), $7\times$ savings on
    $\omegaproj$ parameters, enables higher scale within the same
    VRAM budget.
  \item \textbf{RT core inference.} Speculative. Would allow a
    $5\times$ wider readout within the same speed budget.
  \item \textbf{Hybrid SSM + attention.} A 3:1 SSM-to-attention ratio
    in the layer stack, borrowing from state-space model research.
    Open experiment; not yet piloted.
  \item \textbf{Int6 quantization-aware training.} Reduces the HBM
    pressure and enables larger effective batch; speculative on bpb
    impact.
\end{enumerate}

The first three are \emph{known} to work in expectation. The last three
are research directions.

% ============================================================
\part{Open Questions}
% ============================================================

\chapter{The Additive Law}

Several Session 11 observations are consistent with an \emph{additive law}
of bpb improvements under fixed data and scale:
\begin{equation}
  \bpb(\text{A + B}) \approx \bpb(\text{A}) + \Delta_B,
\end{equation}
where $\Delta_B$ is the standalone improvement of mechanism B. Examples:
persistent state at s8 improves by 0.016; at s12 we observe 0.015. The
lasso improves by 0.109 at s8; at s12 it improves by 0.092 (smaller, but
consistent).

The additive law has downside risk when mechanisms interact. Heinrich
predicts that adding persistent state on top of learnable-$\omegaproj$
will deliver $-0.005$ to $-0.015$ rather than the additive $-0.016$,
because the expanded substrate (learnable-s8 $\effdim = 314$) has less
to compress.

This conjecture is the subject of the currently-running pilot
\texttt{mut-routed-learnable-persistent-s12-50k}.

\begin{conjecture}[Bounded additivity]
Under fixed scale and data, mechanism combinations are additive in bpb
improvement up to a multiplicative correction that depends on the
overlap of the effective-dimension subspaces the mechanisms address.
\end{conjecture}

\chapter{The Scan Is Dead}

Session 10's $\cka = 1.000$ finding: consecutive substrate states are
identical up to CKA precision. The EMA isn't doing any work; the
signal is entirely in $\omegaproj$ and in the local path.

This raises a question: why keep the scan at all? An ablation removing
the scan entirely changes bpb by $< 0.01$.

Two hypotheses, neither yet settled:
\begin{enumerate}
  \item The scan is infrastructure. It exists to thread the
    $\omegaproj$-modulated input into a sequence-length-independent
    representation. Even though the values don't change, the position
    in the substrate does --- and the readout reads at the correct
    position.
  \item The scan is literally dead weight. Remove it and the same
    bpb is reachable by a purely local model. This would be a major
    architectural simplification.
\end{enumerate}

The empirical test --- run learnable-s8 with scan removed ---
has not been conducted as of this writing. It is the cheapest
high-leverage experiment remaining in the Session 11 queue.

\chapter{Forecasting}

Chronohorn's \texttt{record\_result} fits a power-law saturation curve
to the probe sequence, producing a forecast asymptote. Session 10's
diagnosis: \emph{the forecast is biased $+0.18$ bpb optimistic}. Probe
sequences extrapolated to infinity predict 1.6 bpb when the actual final
is 1.78.

Two contributors:
\begin{itemize}
  \item The probe-final gap: a 16-batch probe's bpb is $\sim 0.1$ higher
    than a 32-batch final in training regimes where evaluation has
    warmed up. Probes are a leading indicator, not the final.
  \item The asymptote fit is pulled low by the rapid early descent.
    A two-parameter power law doesn't have the flexibility to fit both
    the steep early drop and the slow late flatten.
\end{itemize}

Fix candidates: a three-parameter fit; weighting the fit toward the
tail; using a longer warmup before probes count. None implemented at
time of writing.

\chapter{The Frontier Question}

The top of the frontier table currently shows
\texttt{acb-byte-s1024-10k} at 0.866 bpb. The ``byte'' prefix and
the naming scheme suggest this is a byte-level model; the config
JSON's \texttt{vocab\_size} field is null, making it impossible to
verify from the DB alone. If the run is actually tokenized (sp1024) and
has been misclassified into the byte frontier, the apparent best is
off by $2.4\times$.

This is known as the \emph{frontier class} bug. It is unresolved as of
this book. Suggested fix: require \texttt{vocab\_size} in the config
blob, and enforce class separation in the frontier query via
\texttt{WHERE vocab\_size = 256} for byte-level, or an explicit
\texttt{family} column.

% ============================================================
\appendix
\part{Appendices}
% ============================================================

\chapter{Command Reference}
\label{ap:commands}

\section{Runtime}

\begin{lstlisting}[language=bash,caption={Start the unified runtime}]
# Monitor-only (no dispatch, safe default)
chronohorn runtime --port 7071 --no-dispatch

# With a manifest (will dispatch from it)
chronohorn runtime --port 7070 --manifest manifests/my_scan.jsonl

# With auto-deepen
chronohorn runtime --manifest manifests/scan.jsonl --auto-deepen --max-steps 50000
\end{lstlisting}

\section{MCP stdio}

\begin{lstlisting}[language=bash]
chronohorn mcp  # speaks MCP protocol on stdin/stdout
\end{lstlisting}

Configuration in \texttt{.mcp.json}:
\begin{lstlisting}[language=bash]
{
  "mcpServers": {
    "chronohorn": {
      "command": "python",
      "args": ["-m", "chronohorn.mcp_transport"],
      "env": { "PYTHONPATH": "python" }
    }
  }
}
\end{lstlisting}

\section{Fleet}

\begin{lstlisting}[language=bash]
# Register a manually-launched run
chronohorn_register_run name=my-run host=slop-01 steps=50000 seed=42

# Sync fleet: pull results, probe hosts
chronohorn_fleet_sync hosts=["slop-01","slop-02"]

# Clean poisoned DB
chronohorn_reset
\end{lstlisting}

\section{Data Provisioning}

\begin{lstlisting}[language=bash]
chronohorn data provision                # bytes
chronohorn data provision --variant sp4096
chronohorn data provision --verify
\end{lstlisting}

\section{Common Trainer Flags (causal-bank)}

\begin{itemize}
  \item \texttt{--vocab-size 256}: byte level.
  \item \texttt{--scale 8.0}: width scale.
  \item \texttt{--steps 20000}: training steps.
  \item \texttt{--seq-len 1024}.
  \item \texttt{--batch-size 32}: or let auto-scale pick.
  \item \texttt{--substrate-mode frozen}.
  \item \texttt{--adaptive-substrate}: turn on adaptive path.
  \item \texttt{--hrr-omega-init}: HRR Fourier basis init for
    $\omegaproj$'s bias.
  \item \texttt{--triton-scan}: Triton fused scan.
  \item \texttt{--linear-readout-kind routed\_sqrelu\_experts}: routed
    readout.
  \item \texttt{--linear-readout-num-experts 2}.
  \item \texttt{--lasso-rotation}: the lasso.
  \item \texttt{--persistent-state}: persistent across batches.
  \item \texttt{--save-checkpoint-every 5000}: periodic checkpoints.
  \item \texttt{--final-eval-batches 32}: final eval count.
\end{itemize}

\chapter{Result JSON Format}

\begin{lstlisting}[caption={Canonical result JSON shape},basicstyle=\ttfamily\footnotesize]
{
  "model": {
    "test_bpb": 1.785,
    "test_bits_per_token": 1.785,
    "architecture": "causal_bank_routed_experts",
    "params": 17800000
  },
  "config": {
    "train": {
      "steps": 50000,
      "seq_len": 1024,
      "batch_size": 8,
      "learning_rate": 0.001,
      "vocab_size": 256
    }
  },
  "dataset": {
    "test_tokens_per_byte": 1.0,       # bytes: 1.0
    "test_bytes_per_token": 1.0
  },
  "training": {
    "final_eval_batches": 32,
    "performance": {
      "tokens_per_second": 316000,
      "elapsed_sec": 10800,
      "steps_completed": 50000,
      "estimated_sustained_tflops": 10.26
    },
    "probes": [
      {"step": 200,  "bpb": 4.2, "eval_batches": 4},
      {"step": 800,  "bpb": 3.1, "eval_batches": 8},
      {"step": 5000, "bpb": 2.4, "eval_batches": 16},
      {"step": 50000, "bpb": 1.79, "eval_batches": 32}
    ]
  }
}
\end{lstlisting}

Required fields: \texttt{model.test\_bpb} or
\texttt{model.test\_bits\_per\_token} (plus
\texttt{dataset.test\_tokens\_per\_byte}). Recommended:
\texttt{model.architecture} for family detection.
Probe-aware: \texttt{probes[*].eval\_batches} to distinguish monitoring
from claims.

\chapter{Manifest JSONL Format}

One JSON object per line. Fields:

\begin{itemize}
  \item \texttt{name}: unique run identifier.
  \item \texttt{family}: ``causal-bank'' / ``polyhash'' / ``transformer''.
  \item \texttt{architecture}: model architecture string.
  \item \texttt{command}: shell command to run inside the container.
  \item \texttt{hosts}: preferred host list (\texttt{["slop-01",
    "slop-02"]}).
  \item \texttt{image}: docker image (e.g.
    \texttt{pytorch/pytorch:2.8.0-cuda12.8-cudnn9-runtime}).
  \item \texttt{min\_gpu\_mem\_gb}: minimum GPU memory in GB.
  \item \texttt{snapshot\_paths}: source dirs to snapshot into the
    container.
  \item \texttt{source\_dir}: local root for snapshots.
  \item \texttt{remote\_cwd\_rel}: working directory inside the
    snapshot.
  \item \texttt{workload\_kind}: ``training.frontier'' /
    ``training.mutation'' / ``forensics''.
  \item \texttt{scale}, \texttt{steps}, \texttt{seq\_len}, \texttt{seed},
    \texttt{batch\_size}: for predicted-cost calculation.
\end{itemize}

Comments (lines starting with \texttt{\#}) are allowed and ignored.

\chapter{Database Schema (selected columns)}

\begin{lstlisting}[caption={results},basicstyle=\ttfamily\footnotesize]
CREATE TABLE results (
  name TEXT PRIMARY KEY,
  config_id INTEGER,
  bpb REAL NOT NULL,            -- canonical bpb (always byte-level)
  train_bpb REAL,
  overfit_pct REAL,
  tok_s REAL,
  tflops_s REAL,
  total_tflops REAL,
  wall_sec REAL,
  slope REAL,
  illegal INTEGER,              -- family-adapter-detected illegal config
  json_archive TEXT,            -- path to original JSON
  created_at REAL,
  family TEXT,
  steps INTEGER,
  seq_len INTEGER,
  params INTEGER,
  int6_mb REAL,                 -- int6-quantized size
  eval_batches INTEGER,         -- distinguishes probe from final
  accelerator_arch TEXT
);
\end{lstlisting}

\begin{lstlisting}[caption={probes},basicstyle=\ttfamily\footnotesize]
CREATE TABLE probes (
  id INTEGER PRIMARY KEY AUTOINCREMENT,
  name TEXT NOT NULL,
  step INTEGER NOT NULL,
  bpb REAL NOT NULL,
  loss REAL,
  elapsed_sec REAL,
  train_elapsed_sec REAL,
  eval_batches INTEGER,
  UNIQUE(name, step)
);
\end{lstlisting}

\begin{lstlisting}[caption={jobs},basicstyle=\ttfamily\footnotesize]
CREATE TABLE jobs (
  name TEXT PRIMARY KEY,
  manifest TEXT,
  config_id INTEGER,
  state TEXT,                   -- pending/dispatched/running/completed/failed
  launched_at REAL,
  completed_at REAL,
  failure_reason TEXT,
  failure_log TEXT,
  host TEXT,
  remote_run TEXT,              -- /tmp/chronohorn-runs/<name>
  runtime_namespace TEXT,
  runtime_job_name TEXT,
  runtime_pod_name TEXT,
  runtime_node_name TEXT,
  executor_kind TEXT,
  executor_name TEXT,
  launcher TEXT,
  container TEXT,
  checkpoint_path TEXT,
  checkpoint_host TEXT,
  steps INTEGER, seed INTEGER, lr REAL, batch_size INTEGER,
  command TEXT,
  job_json TEXT                 -- full manifest row
);
\end{lstlisting}

\chapter{Glossary}

\begin{description}
  \item[Adaptive substrate] Substrate where $\omegaproj$ produces
    input-dependent modulation of the HRR basis.
  \item[bpb] Bits per byte; the canonical measurement. Defined in
    Equation \ref{eq:bpb}.
  \item[bpt] Bits per token; model-internal metric, must be converted
    to bpb before comparison.
  \item[Catch-up] Pipeline path that rescues results stranded on hosts
    by re-invoking \texttt{pull\_remote\_result}.
  \item[Causal bank] The family of models anchored by Equation
    \ref{eq:substrate}.
  \item[Chronohorn] The runtime: DB, fleet, MCP, dashboard.
  \item[CKA] Centered Kernel Alignment; a similarity measure between
    two activation subspaces.
  \item[ConR$^2$] Content $R^2$; linear regression of substrate on
    current-byte one-hot.
  \item[Decepticons] The model library.
  \item[Drain] The loop that dispatches, pulls, and reconciles.
  \item[EffDim] Effective dimensionality.
  \item[FineWeb 10B] The test corpus; $\sim$10 GB of filtered web text.
  \item[Frontier] The admissible best $k$ results by bpb, for a
    specified evaluation class.
  \item[Ghost] A dispatched/running job the fleet can no longer see;
    older than 6 hours and absent from the k8s pod list.
  \item[Heinrich] The forensic suite.
  \item[HRR] Holographic Reduced Representation; elementwise
    multiplication in a Fourier basis.
  \item[Lasso] The $2 \times 2$ noncommutative rotation primitive.
  \item[MCP] Model Context Protocol; the tool surface.
  \item[Marginal] Incremental bpb per TFLOP; the efficiency metric
    used in the dashboard's eff pane.
  \item[Participation ratio] An alternative effective-rank measure;
    $\mathrm{PR}(H) = (\sum \lambda_i)^2 / \sum \lambda_i^2$ (same
    formula as EffDim, applied to different covariance).
  \item[Persistent state] Substrate state carried across batches,
    enabled by \texttt{--persistent-state}.
  \item[Probe] Short evaluation (2--16 batches) during training, for
    monitoring.
  \item[PosR$^2$] Position $R^2$; linear regression of substrate on
    one-hot position.
  \item[Routed experts] Readout with a softmax router over $N$ experts.
  \item[Scale $s$] Width scale; $d \approx 40 s$, $M \approx 256 s$.
  \item[Sharts] The shared NAS volume at \texttt{/Volumes/Sharts}
    holding checkpoints and MRI artifacts.
  \item[sp1024] Sentencepiece tokenizer with vocab 1024; historical
    baseline family.
  \item[$\omegaproj$] The learned projection from embedding to
    modulation vector; shape $(M, d)$.
\end{description}

\chapter{Commit Hashes Referenced}

\begin{itemize}
  \item \texttt{00807bc} --- Session 9 paper.
  \item \texttt{1a76540} --- Session 9 features (lasso, adaptive,
    byte-level, fleet fixes).
  \item \texttt{6d6ad04} --- Session 8 paper.
  \item \texttt{30c4083} --- Hash memory pilot manifests + CLI flags.
  \item \texttt{50b5f3d} --- Hash memory plumbing fix (silent bypass).
  \item \texttt{62dbce2} --- Added \texttt{omega\_lr\_mult} to
    optimizer param groups.
  \item \texttt{3eaf0aa} --- Gradient-sign fixtures for adaptive
    substrate plumbing.
  \item \texttt{261785a} --- Checkpoint pull enumerates periodic
    files.
  \item \texttt{575acb7}, \texttt{b80975e} --- Per-job path in
    \texttt{chronohorn\_fleet\_pull}.
  \item \texttt{3bc5743} --- Preflight gradient-sign check;
    \texttt{init\_signature\_sha256}; periodic save-checkpoint flag.
  \item \texttt{f52865e} --- Heinrich loader collapse fix.
  \item \texttt{98e5f78} --- Catch-up path for completed jobs
    missing Sharts.
  \item \texttt{7392443} --- Widened catch-up to state='failed'.
  \item \texttt{4acb6ce} --- P0 sweep: ghost rescue, tag hygiene,
    async catch-up.
  \item \texttt{8c28a35} --- Dashboard UI fixes.
  \item \texttt{d4e3f49} --- Session 11 ablations manifest.
\end{itemize}

\backmatter

\chapter*{Afterword: The Collaboration}
\addcontentsline{toc}{chapter}{Afterword}

This book is the product of a collaboration between an individual
researcher (asuramaya) and a language model collaborator (Claude Opus
4.7 with 1M context). The division of labor is worth naming because it
is specific and repeatable:

\begin{itemize}
  \item \emph{Direction and judgement}: researcher. Which mechanism to
    try next, which result is important, which bug is worth chasing
    and which is tolerable --- these are always decisions made by a
    human weighing the state of the project against a longer-term
    goal.
  \item \emph{Execution bandwidth}: model. Writing manifests,
    authoring the drain code, debugging the auto-pull pipeline,
    drafting the chapters of this book --- all done by the model
    under supervision.
  \item \emph{Bookkeeping and memory}: shared. The researcher keeps
    the big picture; the model keeps the memory system (a directory
    of markdown files accessed on every conversation) and the
    database (the single source of truth for results).
\end{itemize}

The project's ratio of research-questions-considered to
experiments-conducted-successfully is roughly an order of magnitude
higher than what the researcher could sustain solo. The price is that
every observation in this book has a specific provenance: a commit,
a Heinrich report, or a DB row. Without that provenance, the
collaboration couldn't verify its own claims, and would drift.

The goal of Chronohorn is a language model. The side effect is a
methodology for conducting research at this scale with one person and
one model, at a pace that was previously available only to teams.

\end{document}
